\chapter{Replicación de sistemas complejos naturales}
\label{chap:naturales}

El autómata celular neuronal tiene la capacidad de replicar sistemas bidimensionales con una regla muy bien definida. La regla obtenida no pudo replicar todos los escenarios en todos los experimentos; sin embargo, sí conserva una fuerte adherencia al comportamiento original. 

Hasta el momento, solo se experimentó con simulaciones en donde tenemos un control absoluto de todos los escenarios, pues al tener una regla bien definida y, más aún, un sistema del cual obtener los datos de entrenamiento, tenemos la referencia exacta de cómo la evolución debiera ser. Pero en la vida real no es posible obtener todas estas evoluciones, y los siguientes experimentos buscan dar paso al objetivo central de esta investigación: brindar a los autómatas celulares la capacidad de aprender reglas del entorno con la referencia evolutiva y no con la regla del sistema. 

Dicho esto, en esta parte se siguen tres experimentos principales que buscan dar propósito a este sistema: la simulación de la propagación de un incendio, la simulación de tormentas y el crecimiento de un moho mucilaginoso.

